% The hardware figure for \S7. Kept in its own file so it can be redrawn without
% touching main.tex.
% Encoding: vermillion is the control plane, green marks Rio's two additions (the dashed
% second bank in each double-buffered table, and the survivor link), and the two packet
% walks are black. The walks had a color each until the two additions needed one; they do
% not cross, so one color carries both, and dropping the blue/green pair also drops the
% figure's least separable hues.
% The PE panels are stacked to match tree depth, so the commit's taps read top-down.
% The survivor link is deliberately small and sits *below* the pipeline row rather than
% in it, per Zhiyuan: Rio's PE is meant to look like a usual PIFO-tree PE, and a fifth
% full-height block made the link look like a component on the scale of the rank logic.
% It is a block in the dequeue path's order rather than an arrow, because he asked that
% no path fork or double back; and the dequeue meets it *before* the PIFOs, since the
% link is what decides which PIFO the pop reads.
% Each walk is one route shaped like an "n", labelled at its tail, and the two labels
% line up. Every block is tall enough to leave a clear band across its bottom, and the
% routes are drawn last so they run *in front of* the blocks, through that band, rather
% than showing only in the gaps between them.
% Enqueue rises a quarter of the way into classify and drops a third of the way into
% the PIFOs; dequeue rises two thirds of the way into them, which is where the survivor
% link hangs, and drops three quarters of the way into next hop. The two outer turns are
% off-centre so neither arrowhead lands on the survivor link's label.
% Both descending stems run well past the block edge they cut, so the arrowhead is not
% read as part of that edge.
% The commit taps are drawn after the PE panels so nothing can clip their arrowheads.
\definecolor{ctlcol}{HTML}{D55E00}
\definecolor{newcol}{HTML}{009E73}
\begin{figure}[t]
\centering
\begin{tikzpicture}[
  x=1cm, y=1cm,
  >={Stealth[length=4pt]},
  pebox/.style={draw=black!45, rounded corners=2.5pt, line width=0.6pt},
  unit/.style={draw=black!45, rounded corners=1.5pt, line width=0.5pt},
  txn/.style={unit, fill=black!8},
  bar/.style={rounded corners=0.6pt, line width=0.5pt},
  ul/.style={font=\scriptsize, inner sep=1pt, anchor=north, align=center},
  plab/.style={font=\scriptsize\itshape, text=black!55, inner sep=1pt},
  tl/.style={font=\scriptsize, inner sep=1pt},
  ctl/.style={draw=ctlcol, line width=0.75pt},
  new/.style={draw=newcol, line width=0.6pt},
  walk/.style={draw=black!75, line width=0.8pt},
]

% ===================== transaction controller =====================
\draw[pebox] (0,4.30) rectangle (8.2,5.25);
\node[font=\scriptsize\bfseries, anchor=west, inner sep=1pt] at (0.18,4.70)
  {transaction controller};
% the instruction queue: the leading four cells are one transaction
\foreach \i in {0,...,8}{%
  \pgfmathsetmacro{\cx}{3.50+0.40*\i}%
  \ifnum\i<4
    \draw[unit, fill=ctlcol!18] (\cx,4.56) rectangle (\cx+0.40,4.84);
  \else
    \draw[unit] (\cx,4.56) rectangle (\cx+0.40,4.84);
  \fi}
\draw[decorate, decoration={brace, amplitude=3pt, raise=1pt},
      draw=black!55, line width=0.5pt] (3.50,4.86) -- (5.10,4.86);
\node[tl, anchor=south] at (4.30,5.00) {one transaction};
\draw[->, draw=black!55, line width=0.5pt] (7.70,4.70) -- (7.16,4.70);

% ===================== the commit trunk, delay-aligned =====================
% One trunk out of the queue's head, tapping each level at its own instant.
\draw[ctl] (3.70,4.56) -- (3.70,4.10) -- (0.62,4.10) -- (0.62,-0.07);
\node[tl, text=ctlcol, rotate=90, anchor=center] at (0.32,1.90) {delay-aligned commit};

% ===================== the processing elements =====================
\draw[pebox] (1.05,1.36) rectangle (8.2,3.75);
\draw[pebox] (1.05,0.54) rectangle (8.2,1.00);
\draw[pebox] (1.05,-0.30) rectangle (8.2,0.16);
\node[plab, anchor=south west] at (1.05,3.78) {PE, level 0};
\node[plab, anchor=south west] at (1.05,1.03) {PE, level 1};
\node[plab, anchor=south west] at (1.05,0.19) {PE, level 2};

% The deeper levels repeat level 0's pipeline row, with their contents left out.
% Same styles as level 0: the gray is a legend color, so it cannot vary by level.
\foreach \yb/\yt in {0.63/0.91, -0.21/0.07}{%
  \draw[txn]  (1.18,\yb) rectangle (2.80,\yt);
  \draw[unit] (3.00,\yb) rectangle (4.50,\yt);
  \draw[unit] (4.70,\yb) rectangle (6.24,\yt);
  \draw[txn]  (6.44,\yb) rectangle (8.08,\yt);}

% ===================== level 0, drawn out =====================
\draw[txn]  (1.18,2.44) rectangle (2.80,3.66);
\draw[unit] (3.00,2.44) rectangle (4.50,3.66);
\draw[unit] (4.70,2.44) rectangle (6.24,3.66);
\draw[txn]  (6.44,2.44) rectangle (8.08,3.66);
\node[ul] at (1.99,3.64) {classify};
\node[ul] at (3.75,3.64) {rank logic};
\node[ul] at (5.47,3.64) {PIFOs};
\node[ul] at (7.26,3.64) {next hop};

% Rio's first addition: a second bank for each match-action table, drawn dashed in
% green under the active gray one. Both mappers take the commit, which Zhiyuan confirmed.
\draw[bar, draw=black!55, fill=black!16] (1.30,3.14) rectangle (2.68,3.34);
\draw[new, bar, dashed]                  (1.30,2.84) rectangle (2.68,3.04);
\draw[bar, draw=black!55, fill=black!16] (6.56,3.14) rectangle (7.96,3.34);
\draw[new, bar, dashed]                  (6.56,2.84) rectangle (7.96,3.04);

% rank-assignment logic is not double-buffered: its writes apply as they arrive
\draw[draw=black!50, fill=black!3, line width=0.5pt]
  (3.44,3.21) -- (4.06,3.06) -- (4.06,2.94) -- (3.44,2.79) -- cycle;

% two logical PIFOs in one PE, which is what spares the survivor a level of its own
\draw[bar, draw=black!55] (4.91,2.96) rectangle (5.41,3.22);
\draw[bar, draw=black!55] (5.53,2.96) rectangle (6.03,3.22);

% ===== Rio's second addition: the survivor link =====
% Small, below the row, on the dequeue's way in. A stored pointer the dequeue reads
% before it picks a PIFO to pop. Labelled on its right, so the enqueue's descent does
% not appear to run into the words.
\draw[new, rounded corners=1.5pt] (5.53,1.96) rectangle (6.63,2.26);
\fill[newcol] (5.93,2.11) circle (1.3pt);
\draw[->, draw=newcol, line width=0.5pt] (5.99,2.11) -- (6.27,2.11);
\node[tl, text=newcol, anchor=west, align=left] at (6.71,2.11) {survivor\\link};

% ===================== the commit's taps, one per level =====================
\draw[ctl,->] (0.62,3.30) -- (1.05,3.30);
\draw[ctl,->] (0.62,0.77) -- (1.05,0.77);
\draw[ctl,->] (0.62,-0.07) -- (1.05,-0.07);
\node[tl, text=ctlcol, anchor=south] at (0.88,3.37) {\(t_0\)};
\node[tl, text=ctlcol, anchor=south] at (0.84,0.84) {\(t_1\)};
\node[tl, text=ctlcol, anchor=south] at (0.84,0.00) {\(t_2\)};

% ===== the two walks, drawn last so they run in front of the blocks =====
% Enqueue: up a quarter into classify, right along the band, down a third into the PIFOs.
\draw[walk,->] (1.585,1.74) -- (1.585,2.64) -- (5.213,2.64) -- (5.213,1.96);
% Dequeue: in over the survivor link, up two thirds into the PIFOs, right along the
% band, down three quarters into next hop. One unbroken route, running in front of the
% link box as it does in front of every other block.
\draw[walk,->] (5.727,1.74) -- (5.727,2.64) -- (7.67,2.64) -- (7.67,1.96);
\node[tl, anchor=north] at (1.585,1.70) {enqueue};
\node[tl, anchor=north] at (5.727,1.70) {dequeue};

\end{tikzpicture}
\caption{A PIFO tree substrate with Rio's two additions (green): a second bank for each match--action table, and a designated-survivor link. The gray blocks are double-buffered, and a commit swaps the banks when it reaches each level at some instant $t_i$.}
\label{fig:hardware}
\end{figure}
